
\subsubsection{Models and Dataset}

\begin{table*}[htbp]
\centering
\resizebox{\dimexpr 2\columnwidth+\columnsep\relax}{!}{%
\begin{tabular}{lllll}
\toprule
Model & Parameters & Layers & Hidden Dim & Attention \\
\midrule
BERT-base-uncased & 110M & 12 & 768 & Bidirectional \\
GPT-2 & 124M & 12 & 768 & Causal \\
\bottomrule
\end{tabular}
}
\end{table*}

\textbf{Dataset:} SST-2 (GLUE benchmark) with 5\% injected label noise.
\begin{itemize}
\item Training samples: 67,349
\item Validation samples: 872
\item Mislabeled examples: 3,367 (5\%)
\end{itemize}

\subsubsection{Experimental Configuration}

\begin{table}[htbp]
\centering
\resizebox{\columnwidth}{!}{%
\begin{tabular}{ll}
\toprule
Parameter & Value \\
\midrule
Seeds & 42, 123 \\
Compute Budgets & 64, 256, 1024 (projection dimension) \\
Train Samples (attribution) & 500 \\
Query Samples & 100 \\
Statistical Threshold & $\alpha$ = 0.05 \\
Fine-tuning Epochs & 3 \\
Learning Rate & 2e-5 \\
Batch Size & 32 \\
Optimizer & AdamW \\
Device & CUDA \\
\bottomrule
\end{tabular}
}
\end{table}

\subsubsection{Statistical Analysis}

All experiments use 2 random seeds (42, 123). Mean $\pm$ std is reported. Statistical comparisons use paired t-tests with p<0.05 threshold. This sample size provides directional evidence but limited statistical power; the authors note that 5+ seeds would be needed for definitive confirmation of P1/P2.
